\documentclass[12pt]{article}

\usepackage[margin=1in]{geometry}
\usepackage{amsmath,amssymb}
\usepackage{enumitem}

\newcommand\E{\mathbb{E}}

\setlist{nosep}

\begin{document}

\noindent
\textbf{Name:} \rule{3in}{0.4pt}

\vspace{0.5cm}

\begin{center}
    {\Large \textbf{In-Class Activity: Experiments}}\\[0.2cm]
    {\large PS200C: Causal Inference in the Social Sciences}\\
    {\large Spring 2026}
\end{center}

\vspace{0.5cm}

\begin{enumerate}

\item \textbf{From ATE to the Difference in Means}

Suppose treatment $D_i$ is randomly assigned, so that $\{Y_i(1), Y_i(0)\} \perp\!\!\!\perp D_i$.

Starting from the ATE, fill in the missing expressions and justify each step.

\begin{align*}
\tau_{\text{ATE}} &= \E[Y_i(1)] - \E[Y_i(0)] \\[6pt]
    &= \E[Y_i(\rule{.4cm}{0.4pt}) \mid D_i= \rule{.4cm}{0.4pt}] - \E[Y_i(\rule{.4cm}{0.4pt}) \mid D_i=\rule{.4cm}{0.4pt}] \quad \text{ignorability} \\[14pt]
    &= \E[\rule{.4cm}{0.4pt} \mid D_i= 1] - \E[\rule{.4cm}{0.4pt} \mid D_i = 0] \quad \text{consistency} \\[14pt]
    &\approx \bar{Y}_1 - \bar{Y}_0
\end{align*}


\vspace{0.3cm}

\item \textbf{Selection bias and randomization}

Recall from the POM lecture that the difference in means can be decomposed as:
\[
\E[Y_i \mid D_i=1] - \E[Y_i \mid D_i=0] \;=\; \underbrace{\E[Y_i(1) - Y_i(0) \mid D_i=1]}_{\text{ATT}} \;+\; \underbrace{\E[Y_i(0)\mid D_i=1] - \E[Y_i(0)\mid D_i=0]}_{\text{Selection Bias}}
\]

\begin{enumerate}[label=(\alph*)]
    \item In words, what does the selection bias term measure?

    \vspace{2cm}

    \item Why does random assignment of $D_i$ eliminate this term?

    \vspace{2cm}

    \item Once the bias term is zero, the DIM equals the ATT. Why does it also equal the ATE under randomization?

    \vspace{2cm}
\end{enumerate}

\clearpage

\item \textbf{Sharp null and randomization inference}

Consider this table of observed data from a small experiment with 2 treated and 2 control units:

\begin{center}
\begin{tabular}{c c c}
$i$ & $D_i$ & $Y_i$ \\
\hline
1 & 1 & 8 \\
2 & 1 & 4 \\
3 & 0 & 3 \\
4 & 0 & 5 \\
\end{tabular}
\end{center}

\begin{enumerate}[label=(\alph*)]
    \item What is the observed difference in means, $\widehat{\tau}$?

    \vspace{0.8cm}

    \item Under the sharp null $H_0: Y_i(1) = Y_i(0)$ for all $i$, fill in $Y_i(1)$ and $Y_i(0)$ for every unit:

    \begin{center}
    {\renewcommand{\arraystretch}{1.6}
    \begin{tabular}{c c c c c}
    $i$ & $D_i$ & $Y_i$ & $Y_i(1)$ & $Y_i(0)$ \\
    \hline
    1 & 1 & 8 & \rule{1cm}{0.4pt} & \rule{1cm}{0.4pt} \\
    2 & 1 & 4 & \rule{1cm}{0.4pt} & \rule{1cm}{0.4pt} \\
    3 & 0 & 3 & \rule{1cm}{0.4pt} & \rule{1cm}{0.4pt} \\
    4 & 0 & 5 & \rule{1cm}{0.4pt} & \rule{1cm}{0.4pt} \\
    \end{tabular}}
    \end{center}

    \vspace{0.3cm}

    \item There are $\binom{4}{2} = 6$ possible assignments. Fill in $\hat\tau(\omega)$ for each, then compute the two-sided p-value.

    \begin{center}
    {\renewcommand{\arraystretch}{1.6}
    \begin{tabular}{c c c}
    Treated units & $\hat\tau(\omega)$ \\
    \hline
    $\{1,2\}$ & \rule{1.5cm}{0.4pt} \\
    $\{1,3\}$ & \rule{1.5cm}{0.4pt} \\
    $\{1,4\}$ & \rule{1.5cm}{0.4pt} \\
    $\{2,3\}$ & \rule{1.5cm}{0.4pt} \\
    $\{2,4\}$ & \rule{1.5cm}{0.4pt} \\
    $\{3,4\}$ & \rule{1.5cm}{0.4pt} \\
    \end{tabular}}
    \end{center}
	
	\vspace{0.5cm}

    Two-sided p-value: \rule{3cm}{0.4pt}

\end{enumerate}

%\vspace{1cm}
%
%\item \textbf{BONUS only: Standard errors}
%
%We discussed the Neyman standard error for the DIM estimator:
%\[
%\widehat{SE} = \sqrt{\frac{\hat\sigma^2_{Y|D=1}}{N_1} + \frac{\hat\sigma^2_{Y|D=0}}{N_0}}
%\]
%
%\begin{enumerate}[label=(\alph*)]
%    \item This SE is exact for which quantity: SATE or PATE?
%
%    \vspace{0.8cm}
%
%    \item For the other quantity, is this SE conservative or anti-conservative? Why, intuitively?
%
%    \vspace{1.5cm}
%\end{enumerate}


\end{enumerate}

\end{document}
